
\section{Related Work}

There are two common paradigms for transfer learning for natural language processing (NLP) tasks, \textit{resource-based} transfer and \textit{model-based} transfer. Resource-based transfer utilizes additional linguistic annotations as weak supervision for transfer learning, such as cross-lingual dictionaries \citep{zirikly2cross}, corpora \citep{wang2013cross}, and word alignments \citep{yarowsky2001inducing}. Resource-based methods demonstrate considerable success in cross-lingual transfer, but are quite sensitive to the scale and quality of the additional resources. Resource-based transfer is mostly limited to cross-lingual transfer in previous works, and there is not extensive research on extending resource-based methods to cross-domain and cross-application settings.

Model-based transfer, on the other hand, does not require additional resources. Model-based transfer exploits the similarity and relatedness between the source task and the target task by adaptively modifying the model architectures, training algorithms, or feature representation. For example, \cite{ando2005framework} proposed a transfer learning framework that shares structural parameters across multiple tasks, and improve the performance on various tasks including NER; \cite{collobert2011natural} presented a task-independent convolutional neural network and employed joint training to transfer knowledge from NER and POS tagging to chunking; \cite{pengimproving} studied transfer learning between named entity recognition and word segmentation in Chinese based on recurrent neural networks. Cross-domain transfer, or domain adaptation, is also a well-studied branch of model-based transfer in NLP. Techniques in cross-domain transfer include the design of robust feature representations \citep{schnabel2014flors}, co-training \citep{chen2011co}, hierarchical Bayesian prior \citep{finkel2009hierarchical}, and canonical component analysis \citep{kim2015new}.

While our approach falls into the paradigm of model-based transfer, in contrast to the above methods, our method focuses on exploiting the generality of deep recurrent neural networks and is applicable to transfer between domains, applications, and languages.

Our work builds on previous work on sequence tagging based on deep neural networks. \cite{collobert2011natural} develop end-to-end neural networks for sequence tagging without hand-engineered features. Later architectures based on different combinations of convolutional networks and recurrent networks have achieved state-of-the-art results on many tasks \citep{collobert2011natural,huang2015bidirectional,chiu2015named,lample2016neural,ma2016end}. These models demonstrate significant generality since they can be applied to multiple applications in multiple languages with a unified network architecture and without task-specific feature extraction.

% Our work is built upon these advances to further exploit the generalizability for transfer learning.


% On one hand, extensive approaches have been developed for transfer learning based on additional linguistic resources, such as cross-lingual dictionaries, corpora, or word alignments.


% Ando and Zhang \shortcite{ando2005framework} proposed a multi-task joint training framework that shares structural parameters among multiple tasks, and improved the performance on various tasks including NER. Collobert et al. \shortcite{collobert2011natural} presented a task independent convolutional network and employed multi-task joint training to improve the performance of chunking. However, there is still a gap between these multi-task approaches and the state-of-the-art results on individual tasks. Furthermore, it is unclear whether these approaches can be effective in a cross-lingual setting.

% Multilingual resources were extensively used for
% cross-lingual sequence tagging through various ways, such as
% cross-lingual feature extraction \cite{darwish2013named}, text
% categorization \cite{virga2003transliteration}, and Bayesian parallel
% data prediction \cite{snyder2008unsupervised}. Parallel corpora and
% word alignments are also used for training cross-lingual distributed
% word representations
% \cite{kiros2014multiplicative,gouws2014bilbowa,zhou2015learning}. Unlike
% these approaches, our method mainly focuses on using morphological
% similarity for cross-lingual joint training.

% Several neural architectures based on recurrent networks were proposed
% for sequence tagging. Huang et al. \shortcite{huang2015bidirectional}
% used word-level Long Short-Term Memory (LSTM) units based on
% handcrafted features; dos Santos et al. \shortcite{dos2015boosting}
% employed convolutional layers on both character and word levels; Chiu
% and Nichols \shortcite{chiu2015named} applied convolutional layers on
% the character level and LSTM units on the word level; Gillick et al. \shortcite{gillick2015multilingual} employed a sequence-to-sequence LSTM with a novel tagging scheme. We show that our
% architecture gives better performance experimentally than these
% approaches in Section \ref{sec:exp}.

% Most similar to our work is the recent approach independently
% developed by Lample et al. \shortcite{lample2016neural} (published two
% weeks before our submission), which employs LSTM on both character and
% word levels. However, there are several crucial differences. First, we
% study cross-lingual joint training and show improvement over their
% approach in various cases. Second, while they mainly focus on NER, we
% generalize our model to other sequence tagging tasks, and also
% demonstrate the effectiveness of multi-task joint training. There are
% also differences in the technical aspect, such as the
% cost-sensitive loss function and gated recurrent units used in our
% work.
